\begin{exercises}
  \item
    转置以下各矩阵。
    \begin{exparts*}
      \partsitem \( \begin{mat}[r]
                 2  &1  \\
                 3  &1
               \end{mat}  \)
      \partsitem \( \begin{mat}[r]
                 2  &1  \\
                 1  &3
               \end{mat}  \)
      \partsitem \( \begin{mat}[r]
                 1  &4  &3 \\
                 6  &7  &8
               \end{mat}  \)
      \partsitem \( \colvec[r]{0 \\ 0 \\ 0} \)
      \partsitem \( \rowvec{-1 &-2} \)
    \end{exparts*}
    \begin{answer}
      \begin{exparts*}
        \partsitem \( \begin{mat}[r]
                   2  &3  \\
                   1  &1
                 \end{mat}  \)
        \partsitem \( \begin{mat}[r]
                   2  &1  \\
                   1  &3
                 \end{mat}  \)
        \partsitem \( \begin{mat}[r]
                   1  &6  \\
                   4  &7  \\
                   3  &8
                 \end{mat}  \)
        \partsitem \( \rowvec{0 &0 &0} \)
        \partsitem \( \colvec[r]{-1 \\ -2} \)
      \end{exparts*}
    \end{answer}
  \recommended \item
    判断该向量是否在矩阵的行空间中。
    \begin{exparts*}
       \partsitem \( \begin{mat}[r]
                  2  &1  \\
                  3  &1
                \end{mat}  \),
             \( \rowvec{1 &0} \)
       \partsitem \( \begin{mat}[r]
                  0  &1  &3  \\
                 -1  &0  &1  \\
                 -1  &2  &7
                \end{mat}  \),
             \( \rowvec{1 &1 &1} \)
    \end{exparts*}
    \begin{answer}
       \begin{exparts}
         \partsitem 是的。
           为判断是否存在 $c_1$ 与 $c_2$ 使得
           \( c_1\cdot\rowvec{2 &1}+c_2\cdot\rowvec{3 &1}=\rowvec{1 &0} \)，
           我们求解
           \begin{equation*}
              \begin{linsys}{2}
                 2c_1  &+  &3c_2  &=  &1  \\
                  c_1  &+  &c_2   &=  &0
              \end{linsys}
           \end{equation*}
           解得 \( c_1=-1 \) 与 \( c_2=1 \)。
           因此该向量在行空间中。
         \partsitem 否。
           方程
           \( c_1\rowvec{0 &1 &3}
              +c_2\rowvec{-1 &0 &1}
              +c_3\rowvec{-1 &2 &7}
              =\rowvec{1 &1 &1} \)
           无解。
           \begin{equation*}
             \begin{amat}[r]{3}
               0  &-1  &-1  &1  \\
               1  &0   &2   &1  \\
               3  &1   &7   &1
             \end{amat}
             \grstep{\rho_1\leftrightarrow\rho_2}
             \grstep{-3\rho_1+\rho_3}
             \repeatedgrstep{\rho_2+\rho_3}
             \begin{amat}[r]{3}
               1  &0   &2   &1  \\
               0  &-1  &-1  &1  \\
               0  &0   &0   &-1
             \end{amat}
           \end{equation*}
           因此，该向量不在行空间中。
      \end{exparts}
    \end{answer}
  \recommended \item
    判断该向量是否在列空间中。
    \begin{exparts*}
      \partsitem \( \begin{mat}[r]
                 1  &1  \\
                 1  &1
               \end{mat}  \),
            \( \colvec[r]{1 \\ 3} \)
      \partsitem \( \begin{mat}[r]
                 1  &3  &1 \\
                 2  &0  &4 \\
                 1  &-3 &-3
               \end{mat}  \),
            \( \colvec[r]{1 \\ 0 \\ 0} \)
    \end{exparts*}
    \begin{answer}
       \begin{exparts}
          \partsitem 否。
            为判断是否存在 \( c_1,c_2\in\Re \) 使得
            \begin{equation*}
              c_1\colvec[r]{1 \\ 1}
              +c_2\colvec[r]{1 \\ 1}
              =\colvec[r]{1 \\ -3}
            \end{equation*}
            我们可以对由此得到的线性方程组使用高斯消元法。
            \begin{equation*}
              \begin{linsys}{2}
                 c_1  &+  &c_2  &=  &1  \\
                 c_1  &+  &c_2  &=  &3
              \end{linsys}
              \grstep{-\rho_1+\rho_2}
              \begin{linsys}{2}
                 c_1  &+  &c_2  &=  &1  \\
                      &   &0    &=  &2
              \end{linsys}
            \end{equation*}
            该方程组无解，因此该向量不在列空间中。
          \partsitem 是的。
            由关系式
            \begin{equation*}
              c_1\colvec[r]{1 \\ 2 \\ 1}
              +c_2\colvec[r]{3 \\ 0 \\ -3}
              +c_3\colvec[r]{1 \\ 4 \\ -3}
              =\colvec[r]{1 \\ 0 \\ 0}
            \end{equation*}
            我们得到一个线性方程组，对其使用高斯消元法，
            \begin{equation*}
              \begin{amat}[r]{3}
                1  &3  &1  &1  \\
                2  &0  &4  &0  \\
                1  &-3 &-3 &0
              \end{amat}
              \grstep[-\rho_1+\rho_3]{-2\rho_1+\rho_2}
              \repeatedgrstep{-\rho_2+\rho_3}
              \begin{amat}[r]{3}
                1  &3  &1  &1  \\
                0  &-6 &2  &-2 \\
                0  &0  &-6 &1
              \end{amat}
            \end{equation*}
            即得一个解。
            因此，该向量在列空间中。
% sage: load('gauss_method.sage')
% sage: M = matrix(QQ, [[1,3, 1, 1], [2, 0, 4, 0], [1, -3, -3, 0]])
% sage: M
% [ 1  3  1  1]
% [ 2  0  4  0]
% [ 1 -3 -3  0]
% sage: gauss_method(M)
% [ 1  3  1  1]
% [ 2  0  4  0]
% [ 1 -3 -3  0]
% (' take', -2, 'times row', 1, 'plus row', 2)
% (' take', -1, 'times row', 1, 'plus row', 3)
% [ 1  3  1  1]
% [ 0 -6  2 -2]
% [ 0 -6 -4 -1]
% (' take', -1, 'times row', 2, 'plus row', 3)
% [ 1  3  1  1]
% [ 0 -6  2 -2]
% [ 0  0 -6  1]
      \end{exparts}
    \end{answer}
  \recommended \item
    判断该向量是否在矩阵的列空间中。
    \begin{exparts*}
      \partsitem \( \begin{mat}[r]
                   2  &1  \\
                   2  &5
                 \end{mat} \),~\( \colvec[r]{1 \\ -3}  \)
      \partsitem \( \begin{mat}[r]
                   4  &-8 \\
                   2  &-4
                 \end{mat} \),~\( \colvec[r]{0 \\ 1}  \)
      \partsitem \( \begin{mat}[r]
                   1  &-1  &1  \\
                   1  &1   &-1 \\
                  -1  &-1  &1
                \end{mat} \),~\( \colvec[r]{2 \\ 0 \\ 0}  \)
    \end{exparts*}
    \begin{answer}
      \begin{exparts}
        \partsitem 是的；
          我们问的是是否存在标量 \( c_1 \) 与 \( c_2 \) 使得
          \begin{equation*}
            c_1\colvec[r]{2 \\ 2}+c_2\colvec[r]{1 \\ 5}=\colvec[r]{1 \\ -3}
          \end{equation*}
          由此得到线性方程组
          \begin{equation*}
            \begin{linsys}{2}
              2c_1  &+  &c_2  &=  &1  \\
              2c_1  &+  &5c_2 &=  &-3
            \end{linsys}
            \grstep{-\rho_1+\rho_2}
            \begin{linsys}{2}
              2c_1  &+  &c_2  &=  &1  \\
                    &   &4c_2 &=  &-4
            \end{linsys}
          \end{equation*}
          高斯消元法给出 \( c_2=-1 \) 与 \( c_1=1 \)。
          也就是说，确实存在这样一对标量，因此该向量
          确实在矩阵的列空间中。
        \partsitem 否；
          我们问的是是否存在标量 $c_1$ 与 $c_2$
          使得
          \begin{equation*}
            c_1\colvec[r]{4 \\ 2}+c_2\colvec[r]{-8 \\ -4}=\colvec[r]{0 \\ 1}
          \end{equation*}
          一种做法是考虑由此得到的线性方程组
          \begin{equation*}
            \begin{linsys}{2}
              4c_1  &-  &8c_2  &=  &0  \\
              2c_1  &-  &4c_2  &=  &1
            \end{linsys}
          \end{equation*}
          容易看出它无解。
          另一种做法是注意到
          左边各列的任意线性组合
          其第二个分量都是第一个分量的一半，
          而右边的向量不满足这一条件。
        \partsitem 是的；我们可以直接观察到该向量
          是第一列减去第二列。
          或者，若一时看不出，可以写出各列之间的关系
          \begin{equation*}
            c_1\colvec[r]{1 \\ 1 \\ -1}
             +c_2\colvec[r]{-1 \\ 1 \\ -1}
             +c_3\colvec[r]{1 \\ -1 \\ 1}
             =\colvec[r]{2 \\ 0 \\ 0}
          \end{equation*}
          并考虑由此得到的线性方程组
          \begin{equation*}
            \begin{linsys}{3}
              c_1  &-  &c_2  &+  &c_3  &=  &2  \\
              c_1  &+  &c_2  &-  &c_3  &=  &0  \\
             -c_1  &-  &c_2  &+  &c_3  &=  &0
            \end{linsys}
            \grstep[\rho_1+\rho_3]{-\rho_1+\rho_2}
            \begin{linsys}{3}
              c_1  &-  &c_2  &+  &c_3  &=  &2  \\
                   &   &2c_2 &-  &2c_3 &=  &-2  \\
                   &   &-2c_2&+  &2c_3 &=  &2
            \end{linsys}
            \grstep{\rho_2+\rho_3}
            \begin{linsys}{3}
              c_1  &-  &c_2  &+  &c_3  &=  &2  \\
                   &   &2c_2 &-  &2c_3 &=  &-2  \\
                   &   &     &   &0    &=  &0
            \end{linsys}
          \end{equation*}
          给出了额外的信息（除了解至少存在一个之外）：有无穷多解。
          参数化给出 $c_2=-1+c_3$ 与 $c_1=1$，取 $c_3$ 为
          零即得一个特解 $c_1=1$，$c_2=-1$，
          $c_3=0$（这当然就是开头所作的观察）。
      \end{exparts}
    \end{answer}
  \recommended \item
    求此矩阵行空间的一个基。
    \begin{equation*}
      \begin{mat}[r]
         2  &0  &3  &4  \\
         0  &1  &1  &-1 \\
         3  &1  &0  &2  \\
         1  &0  &-4 &-1
       \end{mat}
    \end{equation*}
    \begin{answer}
     % Using Sage:
     % M = matrix(QQ, [[2,0,3,4], [0,1,1,-1], [3,1,0,2], [1,0,-4,-1]])
     % M1 = M.with_added_multiple_of_row(2,0,-3/2)
     % M2 = M1.with_added_multiple_of_row(3,0,-1/2)
     % M3 = M2.with_added_multiple_of_row(2,1,-1)
     % M4 = M3.with_added_multiple_of_row(3,2,-1)
     常规的高斯消去
     \begin{equation*}
       \begin{mat}[r]
         2  &0  &3 &4   \\
         0  &1  &1  &-1 \\
         3  &1  &0  &2  \\
         1  &0  &-4  &-1
       \end{mat}
       \grstep[-(1/2)\rho_1+\rho_4]{-(3/2)\rho_1+\rho_3}
       \grstep{-\rho_2+\rho_3}
       \repeatedgrstep{-\rho_3+\rho_4}
       \begin{mat}[r]
         2  &0  &3      &4   \\
         0  &1  &1      &-1 \\
         0  &0  &-11/2  &-3  \\
         0  &0  &0      &0
       \end{mat}
     \end{equation*}
     提示了这个基
     $\sequence{\rowvec{2 &0 &3 &4},
         \rowvec{0 &1 &1 &-1},\rowvec{0 &0 &-11/2 &-3}}$。

      另一个或许更方便的做法是先对换两行，
      % From Sage
      % M = matrix(QQ, [[2,0,3,4], [0,1,1,-1], [3,1,0,2], [1,0,-4,-1]])
      % M1 = M.with_swapped_rows(0,3)
      % M2 = M1.with_added_multiple_of_row(2,0,-3)
      % M3 = M2.with_added_multiple_of_row(3,0,-2)
      % M4 = M3.with_added_multiple_of_row(2,1,-1)
      % M5 = M4.with_added_multiple_of_row(3,2,-1)
      \begin{equation*}
        \grstep{\rho_1\swap\rho_4}
        \repeatedgrstep[-2\rho_1+\rho_4]{-3\rho_1+\rho_3}
        \repeatedgrstep{-\rho_2+\rho_3}
        \repeatedgrstep{-\rho_3+\rho_4}
        \begin{mat}[r]
           1  &0  &-4 &-1 \\
           0  &1  &1  &-1 \\
           0  &0  &11 &6  \\
           0  &0  &0  &0
         \end{mat}
      \end{equation*}
      从而得到基
      \( \sequence{\rowvec{1 &0 &-4 &-1},\rowvec{0 &1 &1 &-1},
                   \rowvec{0 &0 &11 &6}} \)。
    \end{answer}
  \recommended \item
    求每个矩阵的秩。
    \begin{exparts*}
      \partsitem \(
        \begin{mat}[r]
          2  &1  &3  \\
          1  &-1 &2  \\
          1  &0  &3
        \end{mat}  \)
      \partsitem \(
        \begin{mat}[r]
          1  &-1 &2  \\
          3  &-3 &6  \\
         -2  &2  &-4
        \end{mat}  \)
      \partsitem \(
        \begin{mat}[r]
          1  &3  &2  \\
          5  &1  &1  \\
          6  &4  &3
        \end{mat}  \)
      \partsitem \(
        \begin{mat}[r]
          0  &0  &0  \\
          0  &0  &0  \\
          0  &0  &0
        \end{mat}  \)
    \end{exparts*}
    \begin{answer}
      \begin{exparts}
         \partsitem 这一化简
           \begin{equation*}
             \mbox{}
             \grstep[-(1/2)\rho_1+\rho_3]{-(1/2)\rho_1+\rho_2}
             \repeatedgrstep{-(1/3)\rho_2+\rho_3}
             \begin{mat}[r]
               2  &1     &3     \\
               0  &-3/2  &1/2   \\
               0  &0     &4/3
             \end{mat}
           \end{equation*}
           表明行秩（从而秩）为三。
         \partsitem 观察各列可见其余各列都是第一列的倍数（观察各行也可看出同样的结论）。
           因此秩为一。

           或者，化简
           \begin{equation*}
             \begin{mat}[r]
               1  &-1  &2  \\
               3  &-3  &6  \\
               -2 &2   &-4
             \end{mat}
             \grstep[2\rho_1+\rho_3]{-3\rho_1+\rho_2}
             \begin{mat}[r]
               1  &-1  &2  \\
               0  &0   &0  \\
               0  &0   &0
             \end{mat}
           \end{equation*}
           也表明同样的事实。
         \partsitem 这一计算
           \begin{equation*}
             \begin{mat}[r]
               1  &3  &2  \\
               5  &1  &1  \\
               6  &4  &3
             \end{mat}
             \grstep[-6\rho_1+\rho_3]{-5\rho_1+\rho_2}
             \repeatedgrstep{-\rho_2+\rho_3}
             \begin{mat}[r]
               1  &3   &2  \\
               0  &-14 &-9 \\
               0  &0   &0
             \end{mat}
           \end{equation*}
           表明秩为二。
         \partsitem 秩为零。
       \end{exparts}
     \end{answer}
  \item 给出此矩阵列空间的一个基。
    给出该矩阵的秩。
    \begin{equation*}
      \begin{mat}
        1 &3 &-1 &2 \\
        2 &1 &1  &0 \\
        0 &1 &1  &4
      \end{mat}
    \end{equation*}
  \begin{answer}
    我们要求这个张成空间的一个基。
    \begin{equation*}
      \spanof{\colvec{1 \\ 2 \\ 0},
              \colvec{3 \\ 1 \\ 1},
              \colvec{-1 \\ 1 \\ 1},
              \colvec{2 \\ 0 \\ 4}}\subseteq\Re^3
    \end{equation*}
    最直接的做法
    是把这些列转置成行，用高斯消元法求出这些行的张成空间的一个基，
    然后再把它们转置回
    列。
    \begin{equation*}
      \begin{mat}
        1 &2 &0 \\
        3 &1 &1 \\
       -1 &1 &1 \\
        2 &0 &4
      \end{mat}
      \grstep[\rho_1+\rho_3 \\ -2\rho_1+\rho_4]{-3\rho_1+\rho_2}
      \grstep[-(4/5)\rho_2+\rho_4]{(3/5)\rho_2+\rho_3}
      \grstep{-2\rho_3+\rho_4}
      \begin{mat}
        1 &2  &0 \\
        0 &-5 &1 \\
        0 &0  &8/5 \\
        0 &0  &0
      \end{mat}
    \end{equation*}
    舍去零向量（它表明初始向量之间存在冗余），得到列空间的这个基。
    \begin{equation*}
      \sequence{
        \colvec{1 \\ 2 \\ 0},
        \colvec{0 \\ -5 \\ 1},
        \colvec{0 \\ 0 \\ 8/5}
        }
    \end{equation*}
    矩阵的秩是其列空间的维数，所以为三。
    （它也等于其行空间的维数。）
  \end{answer}
  \recommended \item
    求以下每个集合的张成空间的一个基。
    \begin{exparts}
      \partsitem \(
        \set{\rowvec{1 &3},
          \rowvec{-1 &3},
          \rowvec{1 &4},
          \rowvec{2 &1}  }\subseteq\matspace_{\nbym{1}{2}} \)
      \partsitem \(
         \set{\colvec[r]{1 \\2 \\1},
           \colvec[r]{3 \\ 1 \\ -1},
           \colvec[r]{1 \\ -3 \\ -3}  }\subseteq\Re^3  \)
      \partsitem \(  \set{1+x,1-x^2,3+2x-x^2}\subseteq\polyspace_3  \)
      \partsitem \( \set{
          \begin{mat}[r]
            1  &0  &1  \\
            3  &1  &-1
          \end{mat},
          \begin{mat}[r]
            1  &0  &3  \\
            2  &1  &4
          \end{mat},
          \begin{mat}[r]
           -1  &0  &-5 \\
           -1  &-1 &-9
          \end{mat}  }  \subseteq\matspace_{\nbym{2}{3}}  \)
    \end{exparts}
    \begin{answer}
      \begin{exparts}
        \partsitem 这一化简
          \begin{equation*}
             \begin{mat}[r]
               1  &3  \\
              -1  &3  \\
               1  &4  \\
               2  &1
             \end{mat}
             \grstep[-\rho_1+\rho_3 \\ -2\rho_1+\rho_4]{\rho_1+\rho_2}
             \grstep[(5/6)\rho_2+\rho_4]{-(1/6)\rho_2+\rho_3}
             \begin{mat}[r]
               1  &3  \\
               0  &6  \\
               0  &0  \\
               0  &0
             \end{mat}
          \end{equation*}
          给出 \( \sequence{\rowvec{1 &3},\rowvec{0 &6}} \)。
        \partsitem 转置并化简
          \begin{equation*}
             \begin{mat}[r]
               1  &2  &1  \\
               3  &1  &-1 \\
               1  &-3 &-3
             \end{mat}
             \grstep[-\rho_1+\rho_3]{-3\rho_1+\rho_2}
             \begin{mat}[r]
               1  &2  &1  \\
               0  &-5 &-4 \\
               0  &-5 &-4
             \end{mat}
             \grstep{-\rho_2+\rho_3}
             \begin{mat}[r]
               1  &2  &1  \\
               0  &-5 &-4 \\
               0  &0  &0
             \end{mat}
          \end{equation*}
          再转置回去，得到这个基。
          \begin{equation*}
            \sequence{\colvec[r]{1 \\ 2 \\ 1},
              \colvec[r]{0 \\ -5 \\ -4}   }
          \end{equation*}
        \partsitem 首先注意外围空间是
          $\polyspace_3$，而不是 $\polyspace_2$。
          然后，把第一个多项式 $1+1\cdot x+0\cdot x^2+0\cdot x^3$
          取作与行向量 $\rowvec{1 &1 &0 &0}$“相同”，等等，
          就得到
          \begin{equation*}
            \begin{mat}[r]
              1  &1  &0  &0 \\
              1  &0  &-1 &0 \\
              3  &2  &-1 &0
            \end{mat}
            \grstep[-3\rho_1+\rho_3]{-\rho_1+\rho_2}
            \repeatedgrstep{-\rho_2+\rho_3}
            \begin{mat}[r]
              1  &1  &0  &0 \\
              0  &-1 &-1 &0 \\
              0  &0  &0  &0
            \end{mat}
          \end{equation*}
          由此得到基 \( \sequence{1+x,-x-x^2} \)。
        \partsitem 这里“相同”的做法给出
          \begin{equation*}
            \begin{mat}[r]
              1  &0  &1  &3  &1  &-1  \\
              1  &0  &3  &2  &1  &4   \\
             -1  &0  &-5  &-1 &-1 &-9
            \end{mat}
            \grstep[\rho_1+\rho_3]{-\rho_1+\rho_2}
            \repeatedgrstep{2\rho_2+\rho_3}
            \begin{mat}[r]
              1  &0  &1  &3  &1  &-1  \\
              0  &0  &2  &-1 &0  &5   \\
              0  &0  &0  &0  &0  &0
            \end{mat}
          \end{equation*}
          从而得到这个基。
          \begin{equation*}
            \sequence{
            \begin{mat}[r]
              1  &0  &1  \\
              3  &1  &-1
            \end{mat},
            \begin{mat}[r]
              0  &0  &2  \\
             -1  &0  &5
            \end{mat}  }
          \end{equation*}
      \end{exparts}
     \end{answer}
  \item 在自然的向量空间中，给出以下每个集合的张成空间的一个基。
    \begin{exparts}
      \partsitem
         $\set{\colvec{1 \\ 1 \\ 3},
              \colvec{-1 \\ 2 \\ 0},
              \colvec{0  \\ 12 \\ 6}}$
      \partsitem $\set{x+x^2, 2-2x, 7, 4+3x+2x^2}$
    \end{exparts}
    \begin{answer}
      \begin{exparts}
        \partsitem
          把这些列转置成行，化为阶梯形
          （然后去掉零行），再转置回列。
          \begin{equation*}
            \begin{mat}
              1 &1  &3 \\
             -1 &2  &0 \\
              0 &12 &6
            \end{mat}
            \grstep{\rho_1+\rho_2}
            \grstep{-4\rho_2+\rho_3}
            \begin{mat}
              1 &1  &3 \\
              0 &3  &3 \\
              0 &0  &-6
            \end{mat}
          \end{equation*}
          该张成空间的一个基如下。
          \begin{equation*}
            \sequence{
              \colvec{1 \\ 1 \\ 3},
              \colvec{0 \\ 3 \\ 3},
              \colvec{0 \\ 0 \\ -6}
          }
          \end{equation*}
       \partsitem
          与前一小区一样，我们把它们看作行，
          以利用高斯消元法已完成的工作。
          \begin{multline*}
            \begin{mat}
              0 &1  &1 \\
              2 &-2 &0 \\
              7 &0 &0 \\
              4 &3  &2
            \end{mat}
            \grstep{\rho_1\leftrightarrow\rho_2}
            \grstep[2\rho_1+\rho_4]{-(7/2)\rho_1+\rho_3}
            \grstep[-7\rho_2+\rho_4]{-7\rho_2+\rho_3}          \\
            \grstep{-(5/7)\rho_3+\rho_4}
            \begin{mat}
              2 &-2  &0 \\
              0 &1 &1 \\
              0 &0  &-7 \\
              0 &0 &0
            \end{mat}
          \end{multline*}
          该集合的张成空间的一个基是
          $\sequence{2-2x, x+x^2, -5x^2}$。
      \end{exparts}
    \end{answer}
  \item
    哪些矩阵的秩为零？
    秩为一的呢？
    \begin{answer}
      只有零矩阵的秩为零。
      秩为一的矩阵仅有如下形式
      \begin{equation*}
        \begin{mat}
           k_1\cdot \rho  \\
            \vdots  \\
           k_m\cdot \rho
        \end{mat}
      \end{equation*}
      其中 \( \rho \) 是某个非零行向量，且诸 \( k_i \) 不全为零。
      （\textit{注}：
      我们不能简单地说所有行都是第一行的倍数，因为第一行可能是零行。
      \textit{再注}：
      上述结论在把“行”换成“列”后同样适用。）
    \end{answer}
